\chapter{相机采集与图像处理}
\setchapterheader{第\thechapter 章\quad 相机采集与图像处理}

\label{chap:camera}
\begin{intro}
区分相机硬件采集、APP 图传和媒体交互输入，处理 RealSense 相机与系统采集服务的占用关系。
\end{intro}
\section{选择图像入口}
标准版 DC\_V93 的参数见第一章。教育版 EDU\_Pro、EDU\_MAX、EDU\_Tailor 等配有算力板的版本，相机连接算力板，开发程序可以通过硬件 SDK 读取图像。D435i 使用\href{https://github.com/IntelRealSense/librealsense}{librealsense}。

相机采集得到的图像可用于识别、定位、导航或应用逻辑。若还需要给机器人的语音交互 agent 提供画面，将采集结果转换为媒体接口要求的 YUYV 数据，再发布外部视频流。
\section{停止已有采集服务}
机器人原有图传服务可能正在使用相机。开始自己的采集程序前，在算力板停止该服务：
\par\noindent\begin{minipage}{\linewidth}
\begin{command}
sudo systemctl stop noetix-video-capture.service
systemctl status noetix-video-capture.service
\end{command}
\end{minipage}\par

\code{stop} 停止本次运行。若后续开发都由自己的程序采集，可以关闭开机启动：
\par\noindent\begin{minipage}{\linewidth}
\begin{command}
sudo systemctl disable noetix-video-capture.service
\end{command}
\end{minipage}\par

完成开发、准备恢复系统图传时，先结束自己的相机程序，再恢复服务：
\par\noindent\begin{minipage}{\linewidth}
\begin{command}
sudo systemctl start noetix-video-capture.service
sudo systemctl enable noetix-video-capture.service
\end{command}
\end{minipage}\par

\code{start} 启动当前服务，\code{enable} 恢复开机自动启动。
